\documentclass[10pt,a4paper]{article}
\usepackage{fontspec}
\usepackage[margin=1.6cm,top=1.4cm,bottom=1.4cm]{geometry}
\usepackage{booktabs,tabularx,xcolor}
\pagestyle{empty}
\setlength{\parindent}{0pt}
\setlength{\parskip}{4pt}
\definecolor{ok}{HTML}{1B7F3B}
\definecolor{ko}{HTML}{B3261E}
\begin{document}

{\Large\bfseries Test contro un gruppo vero della gara 2022: \texttt{smokeonthewatermark}}\\[2pt]
{\small Gruppo \texttt{fakebitinvestigation} — 21/10/2026 — tutti i numeri sono misurati}

\textbf{Il loro schema.} Wavelet Haar a 2 livelli (sottobande LH2, HL2, HH2), spread spectrum con maschera di energia locale, codice correttore LDPC; forza fissa (non calibrata sul WPSNR). Detection: bit 0/1 decodificati, coseno con soglia 0,5377. \textbf{Banco:} 12 immagini (le 3 della gara 2022 + 0000, 0010, \dots, 0080), watermarkate con il loro codice e il loro watermark (WPSNR 56,23–58,16 dB, media 57,60) e con il nostro (62,05 dB).

\medskip
\textbf{Run 1 — il nostro script di attacco contro di loro} (\texttt{attacchi\_fakebitinvestigation.py}, impostazioni di gara: 25 famiglie + 7 zone + greedy a blocchi con la loro detection come oracolo; 12 processi, 29{,}4 min per immagine)

\begin{tabularx}{\linewidth}{@{}l r X@{}}
\toprule
immagine & WPSNR (dB) & attacco riuscito più leggero \\
\midrule
0070 & 47,79 & mediano 7$\times$7, solo nelle zone scelte \\
0010 & 41,30 & blur $\sigma$ 1 + JPEG QF 16, zone \\
0050 & 41,17 & resize 0,263 (cubic), zone \\
0040 & 39,98 & resize 0,325 (cubic), zone \\
0080 & 39,85 & blur $\sigma$ 0,8 + JPEG QF 19 + blur $\sigma$ 0,6, zone \\
0030 & 39,11 & blur $\sigma$ 1 + JPEG QF 19, zone \\
0060 & 38,95 & resize 0,5 (lab) + JPEG QF 13, zone \\
tree & 38,84 & resize 0,375 (cubic), zone \\
0020 & 38,49 & JPEG QF 8 + blur $\sigma$ 0,8, zone \\
rollercoaster & 37,82 & mediano 5$\times$5 + sharpening $\sigma$ 4 $\alpha$ 0,3, zone \\
0000, buildings & — & \textcolor{ko}{non rotte} (famiglie esaurite in $\sim$9 min, $\sim$190 chiamate) \\
\bottomrule
\end{tabularx}

\textcolor{ok}{\textbf{Rotte 10/12}}, WPSNR medio \textbf{40,33 dB} (mediana 39,48; min 37,82; max 47,79). Tempo totale 19,9 min, 0 errori, 463–606 chiamate per immagine rotta. \emph{Verifica indipendente:} ogni immagine salvata, ripassata alla loro detection, dà 0, e il WPSNR ricalcolato col codice dei lab coincide al centesimo con quello del log.

\medskip
\textbf{Run 2 — i loro attacchi contro le nostre immagini} (le funzioni identiche del loro \texttt{attack.py}; il loro strumento è manuale, quindi lo abbiamo emulato in modo generoso: griglia di 154 attacchi e catene — i parametri suggeriti nei loro commenti, allargati, più blur$\to$JPEG come nel loro esempio — e per ogni immagine il riuscito con il WPSNR più alto secondo la nostra detection, $\tau$ = 0,087)

\begin{tabularx}{\linewidth}{@{}l r X r@{}}
\toprule
immagine & WPSNR (dB) & attacco riuscito più leggero & riusciti / 154 \\
\midrule
0070 & 38,74 & JPEG QF 5 + blur $\sigma$ 0,8 & 12 \\
0040 & 36,58 & JPEG QF 7 & 1 \\
0050 & 36,24 & JPEG QF 5 + blur $\sigma$ 0,8 & 10 \\
0010 & 36,09 & AWGN $\sigma$ 0,01 + JPEG QF 5 & 3 \\
0060 & 35,95 & JPEG QF 6 & 3 \\
0000, 0020, 0030, 0080, buildings, rollercoaster, tree & — & \textcolor{ok}{non rotte} & 0 \\
\bottomrule
\end{tabularx}

\textcolor{ko}{\textbf{Rotte 5/12}}, WPSNR medio \textbf{36,72 dB} (min 35,95; max 38,74): tutte con JPEG a qualità 5–7, il nostro punto debole già noto.

\medskip
\textbf{Confronto e lettura onesta}

\begin{tabularx}{\linewidth}{@{}X r r@{}}
\toprule
 & loro, attaccati da noi & noi, attaccati con i loro attacchi \\
\midrule
WPSNR dell'embedding (media) & 57,60 dB & 62,05 dB \\
immagini rotte & 10/12 & 5/12 \\
WPSNR medio degli attacchi riusciti & 40,33 dB & 36,72 dB \\
fascia di robustezza corrispondente & 5 punti (38–41) & 6 punti (35–38) \\
\bottomrule
\end{tabularx}

\begin{itemize}\itemsep1pt
\item \textbf{Il nostro script funziona su uno schema diverso dal nostro}: 10 immagini su 12, quasi tutte sopra i 38 dB. Il greedy a zone conta: le 10 migliori sono tutte localizzate.
\item \textbf{I due run non sono simmetrici.} Il run 1 è il nostro motore completo (ricerca + zone + oracolo, 30 min per immagine); il run 2 è solo una griglia di attacchi globali con le loro funzioni, senza zone. Un loro operatore che sceglie a mano le zone (\texttt{img\_slice}) potrebbe fare meglio: contro il nostro stesso motore, nella simulazione del 19/10, venivamo rotti in 18 immagini su 27 a 36,5 dB di media.
\item \textbf{A WPSNR diversi.} Loro hanno $\sim$4,5 dB di energia in più nell'embedding (3 punti qualità contro 5) e restano comunque più facili da rompere a WPSNR alto.
\item Non misurato: i falsi positivi della loro detection (soglia vicina al valore del caso, 0,5).
\end{itemize}

\end{document}
